\documentclass[11pt]{article}
\usepackage[margin=1in]{geometry}
\usepackage{amsmath,amssymb,amsthm,mathtools}
\usepackage[hidelinks]{hyperref}
\newtheorem{theorem}{Theorem}
\newtheorem{lemma}[theorem]{Lemma}
\newtheorem{corollary}[theorem]{Corollary}
\title{Prime-power terminal orders for nonconstant binomial supports}
\author{Internal arithmetic proof; no priority or external-review claim}
\date{17 November 2026}
\begin{document}
\maketitle
We give an arithmetic mechanism for infinitely many terminal orders in
the ordinary nonconstant binomial problem, including $X(A+BX^3)$.
It proves neither the all-terminal-order assertion for that support nor
Strong Factorial for every univariate binomial.

For integers $N\ge1$, $C\ge0$ and $d\ge1$, put
\[
 P_{N,C,d}(x)=\sum_{k=0}^N\binom Nk\frac{(C+dk)!}{C!}x^k.
                                                               \tag{1}
\]
Write $s_p(v)$ for the base-$p$ digit sum of $v$.
The digit of $C$ at place $L$ means
$\lfloor C/p^L\rfloor\bmod p$, with the units place numbered zero.

\begin{theorem}\label{thm:main}
Let $p$ be prime, $s\ge1$, $r\ge0$, $1\le m<p$, and set
\[
 d=p^s,\qquad q=p^r,\qquad N=mq,\qquad
 h=\frac{p^{r+s}-1}{p-1}.
\]
If the digit $t$ of $C$ at place $r+s$ satisfies $t+m<p$, then
$P_{N,C,d}$ is squarefree over $\mathbb Q_p$, hence over $\mathbb C$.
Its lower $p$-Newton polygon is the single edge
\[
 (0,0)\longrightarrow(mq,mh).
\]
The normalized residual polynomial on this edge, in its lattice
variable $Y$, is
\[
 {\cal R}(Y)=\sum_{j=0}^m\binom mj(t+1)_j(\epsilon Y)^j,
 \qquad \epsilon=(-1)^h\in\mathbb F_p^\times.                 \tag{2}
\]
In particular, when $m=1$, $P_{p^r,C,p^s}$ is irreducible over
$\mathbb Q_p$ and over $\mathbb Q$.
\end{theorem}

\begin{proof}
Let $c_k$ be the coefficient of $x^k$ in (1).
For $k=jq$, $0\le j\le m$, adding $dk=jp^{r+s}$ to $C$
changes precisely its digit $t$ to $t+j$. The hypothesis rules out
a carry. Legendre's formula and the binomial congruence
$\binom{mq}{jq}\equiv\binom mj\pmod p$ therefore give
\[
 v_p(c_{jq})=
 \frac{djq-s_p(C+djq)+s_p(C)}{p-1}=jh.                       \tag{3}
\]
For $k=jq+\ell$, $0<\ell<q$, let $u=v_p(\ell)<r$.
The binomial coefficient $\binom{mq-1}{k-1}$ is a $p$-adic unit:
its upper base-$p$ digits are $m-1$ followed by $r$ digits $p-1$,
and each digit of $k-1$ is at most the corresponding digit.
This digit test follows by taking coefficients in
$(1+z)^a=\prod_i(1+z^{p^i})^{a_i}$ over $\mathbb F_p$.
Consequently
\[
 v_p\binom{mq}{k}=v_p(mq/k)=r-u.
\]
Digit subadditivity, $s_p(dk)=s_p(k)$, and
$s_p(k)=j+s_p(\ell)\le j+(p-1)(r-u)$ give
\begin{align*}
 v_p(c_k)
 &=v_p\binom Nk+
       \frac{dk-s_p(C+dk)+s_p(C)}{p-1}\\
 &\ge \frac{dk}{p-1}+r-u-\frac{s_p(k)}{p-1}\\
 &\ge \frac{dk-j}{p-1}
 >\frac{dk-k/q}{p-1}=\frac{h}{q}k.
\end{align*}
Thus (3) lists exactly the lattice points on the one lower edge.
Since $h\equiv1\pmod p$, $\gcd(h,q)=1$ and its lattice step is $q$.

For an integer $v\ge0$, put
$U_p(v!)=p^{-v_p(v!)}v!$. Splitting the factorial product into
multiples and nonmultiples of $p$, and iterating, yields
\[
 U_p(v!)\equiv(-1)^{v_p(v!)}\prod_i v_i!\pmod p,               \tag{4}
\]
where $v_i$ are its base-$p$ digits. At each iteration the complete
blocks of nonzero residues have product $-1$.
The digit change $t\mapsto t+j$ in (3) and (4) now give
\[
 \overline{c_{jq}/p^{jh}}
   =\binom mj(-1)^{jh}\frac{(t+j)!}{t!},
\]
which is (2). In particular the residual has full degree $m<p$
and constant coefficient one.

The polynomial $\mathcal R$ satisfies the following identity in
$\mathbb F_p[Y]$, obtained by comparing consecutive coefficients:
\[
 \epsilon Y^2{\cal R}''
 +[1-\epsilon(m-t-2)Y]{\cal R}'
 -\epsilon m(t+1){\cal R}=0.                                \tag{5}
\]
If a nonzero root had multiplicity $e\ge2$, then $e\le m<p$.
The lowest Taylor term in (5) would come only from
$\epsilon Y^2{\cal R}''$, with nonzero factor
$\epsilon Y_0^2e(e-1)$, a contradiction. Zero is not a root.
Thus $\mathcal R$ is squarefree.

For completeness, the Newton residual implication used here is the
weighted-initial-product lemma proved in the
\href{PRIME_POWER_TWO_TERM_SFC.tex}{prime-power two-term note}.
At a slope $h/q$, substitute $x=\pi^{-h}Z$, $\pi^q=p$,
normalize by the smallest valuation, and reduce. Initial polynomials
multiply because the residue ring is a field. For factors over
$\mathbb Q_p$, indices attaining the same weighted minimum differ by
multiples of $q$; after removing the lowest monomial their initial
polynomials are polynomials in $Z^q$. A nonconstant squared factor
would therefore contribute a nonconstant squared residual factor on
some edge. Squarefreeness of (2) excludes it.
This argument concerns the lattice variable $Y$, and does not assert
separability after substituting $Y=Z^q$.

When $m=1$, all roots have valuation $-h/q$. Indeed, at any other
valuation either the constant or the top term in (1) is uniquely
least and cannot cancel. A factor of degree $e$ over $\mathbb Q_p$
therefore has a leading-to-constant valuation difference $eh/q$,
an integer. Since $\gcd(h,q)=1$, $q$ divides $e$.
There is no proper nonconstant factor.
\end{proof}

\section{Consecutive moments for lower exponent one}
For integer Gamma shape $\beta\ge1$ and terminal $N\ge2$, put
$\gamma=N+\beta-1$ and
\[
 T_{N,\gamma,d}(x)=\sum_{k=0}^N\binom Nk(\gamma)_{dk}x^k
                =P_{N,N+\beta-2,d}(x).
\]
The lower-one bridge from the
\href{HIGHER_GAP_ADJOINT_RANK.tex}{higher-gap adjoint note}
identifies a common zero of the moments of $X(A+BX^d)$ at orders
$N-1,N$ with a multiple root of this $T$. To check the parameter,
the normalized upper moment uses shape $\gamma+1$ and is
$T+(d/\gamma)xT'$, while the normalized lower moment uses shape
$\gamma$ and is $T-xT'/N$. The determinant is
$-(\gamma+dN)/(\gamma N)\ne0$; also $T(0)=1$.

\begin{corollary}\label{cor:lowerone}
Let $d=p^s$, $s\ge1$, and $N=mp^r\ge2$, $1\le m<p$.
If the digit of $N+\beta-2$ at place $r+s$ is at most $p-m-1$,
the two Gamma moments of $X(A+BX^d)$ at orders $N-1,N$
cannot both vanish, for arbitrary complex $A,B$ defining a nonzero
phase. In particular, for ordinary factorials ($\beta=1$), this
holds for every such $N$.
\end{corollary}
\begin{proof}
Apply Theorem \ref{thm:main} and the bridge.
For $\beta=1$, $C=N-1=mp^r-1<p^{r+s}$, so its digit $t$ is zero.
If $A$ or $B$ vanishes the phase is a nonzero monomial, whose
positive Gamma moments are all nonzero.
\end{proof}
For gap three this includes every terminal order $N=3^r$ or $2\cdot3^r$
with $N\ge2$. The ordinary assertion is an infinite family of
starting indices, not an all-starting-index theorem for
$X(A+BX^3)$.

\section{A separate consequence for arbitrary lower exponent}
\begin{corollary}\label{cor:generala}
Let $a\ge0$, $d=p^s$ with $s\ge1$, and $N=p^r\ge2$.
If the digit of $a$ at place $s$ is at most $p-2$, then
\[
 M_{N,a,d}(x)=\sum_{k=0}^N\binom Nk(aN+dk)!x^k
\]
is irreducible over $\mathbb Q_p$ and over $\mathbb Q$.
Consequently $M_{N-1,a,d}$ and $M_{N,a,d}$ are coprime, proving
the requested consecutive ordinary factorial assertion at orders
$N-1,N$ for every nonzero $X^a(A+BX^d)$.
\end{corollary}
\begin{proof}
Take $m=1$, $C=aN$ in Theorem \ref{thm:main}. Its digit at place
$r+s$ is the digit of $a$ at place $s$.
Multiplication by the nonzero constant $(aN)!$ preserves
irreducibility. A rational polynomial of degree $N-1$ cannot share
a root with an irreducible rational polynomial of degree $N$.
For $AB\ne0$ the moments are $A^j M_{j,a,d}(B/A)$.
The cases $AB=0$ are nonzero monomials.
\end{proof}

Squarefreeness in Theorem \ref{thm:main} for $m>1$ does not by itself
give this arbitrary-lower-exponent consecutive assertion: the factorial
offset changes with the moment order. The degree argument above uses
irreducibility at $m=1$.

\section{Overlap, exact replay and remaining obligation}
The existing prime-power two-term theorem concerns fixed shapes
$\gamma\le d$ for constant-plus-power phases at \emph{every} index.
It does not cover the growing shape $\gamma=N$ here once $N>d$.
The present proof uses the same factorial-unit and Newton-residual
method, specialized to terminal indices with a single nonzero
base-$p$ digit; it does not improve the existing completed
constant-plus-power conclusion. It also overlaps the completed
gap-one, lower-one gap-two, and specified dyadic equal-support
families in their stated specializations. No novelty claim is made.

The companion exact checker compares the actual coefficient Newton
hull and literal normalized residual with (2), checks (5), tests
the lower-one moment identities, and checks irreducibility-degree
coprimality in bounded instances. Its finite run checks the
implementation; all unbounded quantifiers use the proof above.
Carry inputs are rejected as unsupported rather than classified as
counterexamples to Strong Factorial.
The all-terminal-order gap-three theorem, the other gaps and
unrestricted lower exponents remain open.
Independent replay, formal verification, external review and
priority claims remain false.
\end{document}
